\chapter{Variations de suites}
\thispagestyle{fancy}


\section{Sens de variation d'une suite numérique}

\begin{bexemple}
On a représenté ci-dessous le nuage de points des premiers termes d'une suite $\left(u_{n}\right)$ :
\begin{center}
    \includegraphics[scale=0.4]{1.08.00-Img04}
\end{center}
On peut conjecturer que cette suite est croissante. On constate par exemple que $u_{1}<u_{2}$ ou encore $u_{4}<u_{5}$. De manière générale, on peut écrire : $u_{n}<u_{n+1}$
\end{bexemple}

\begin{bdefis}
Soit une suite numérique $\left(u_{n}\right)$.
\begin{itemDef}
\item La \grasInDef{suite $\left(u_{n}\right)$ est croissante} signifie que pour tout entier $n$, on a $u_{n+1} \geqslant u_{n}$.
\item La \grasInDef{suite $\left(u_{n}\right)$ est décroissante} signifie que pour tout entier $n$, on a $u_{n+1} \leqslant u_{n}$.
\end{itemDef}
\end{bdefis}

\begin{bmethodet}{Étudier les variations d'une suite}
\grasInMet{Énoncé:}
\begin{enumMet}
\item Pour tout $n$ de $\mathbb{N}$, on donne la suite $\left(u_{n}\right)$ définie par : $u_{n+1}=u_{n}+2$.\\
Démontrer que la suite $\left(u_{n}\right)$ est croissante.
\item Pour tout $n$ de $\mathbb{N}$, on donne la suite $\left(v_{n}\right)$ définie par : $v_{n}=4 n+4$.\\
Démontrer que la suite $\left(v_{n}\right)$ est croissante.
\end{enumMet}
\grasInMet{Réponse:}
\begin{enumMet}
\item $u_{n+1}-u_{n}=2>0$\\
On en déduit que $\left(u_{n}\right)$ est croissante.
\item On commence par calculer la différence $v_{n+1}-v_{n}$ :\\
On a $: v_{n}=4 n+4$ donc $v_{\mVert{n+1}}=4\mVert{(n+1)}+4=4 n+4+4=4 n+8$
\begin{align*}
v_{n+1}-v_{n}&=4 n+8-(4 n+4) \\
&=4 n+8-4 n-4 \\
&=4
\end{align*}
On étudie ensuite le signe de $v_{n+1}-v_{n}$ :\\
Pour tout $n$ entier $v_{n+1}-v_{n}=4>0$.\\
On en déduit que la suite $\left(v_{n}\right)$ est croissante.
\end{enumMet}
\end{bmethodet}


\section{variation d'une suite arithmétique}

\begin{bprop}
    $U_n$ est une suite arithmétique de raison $r$.
    \begin{itemDef}
    	\item Si $r > 0$ alors la suite $U_n$ est croissante.
    	\item Si $r < 0$ alors la suite $U_n$ est décroissante.
    \end{itemDef}
\end{bprop}

\begin{bdemo}
    $U_{n+1}-U_n=U_n+r-U_n=r$.
    \begin{itemDem}
    	\item Si $r > 0$ alors $U_{n+1}-U_n>0$  et la suite $U_n$ est croissante.
    	\item Si $r < 0$ alors $U_{n+1}-U_n<0$ et la suite $U_n$ est décroissante.
    \end{itemDem}
\end{bdemo}
\begin{bexemple}
    La suite arithmétique $U_n$ définie par $U_n=8-3n$ est décroissante car de raison négative et égale à $-3$.
\end{bexemple}


\section{variation d'une suite géométrique}

\begin{bprop}
    $U_n$ est une suite géométrique de raison $q$ et de premier terme non nul $U_0$.\\
    Pour $U_0>0$ :
    \begin{itemProp}
    	\item Si $q > 1$ alors la suite $U_n$ est croissante.
    	\item Si $0 < q < 1$ alors la suite $U_n$ est décroissante.
    \end{itemProp}
    Pour $U_0<0$ :
    \begin{itemDef}
    	\item Si $q > 1$ alors la suite $U_n$ est décroissante.
    	\item Si $0 < q < 1$ alors la suite $U_n$ est croissante.
    \end{itemDef}
\end{bprop}

\begin{bdemot}{cas où $U_0 > 0$ }
    $U_{n+1}-U_n=U_0\times q^{n+1}-U_0\times q^n=U_0\times q^n\times (q-1)$.
    \begin{itemDem}
    	\item Si $q > 1$ alors $U_{n+1}-U_n>0$ et la suite $U_n$ est croissante.
    	\item Si $0 < q < 1$ alors $U_{n+1}-U_n<0$ et la suite $U_n$ est décroissante.
    \end{itemDem}
\end{bdemot}




